\section{生成规则、候选排序与评价}
\label{sec:evaluation}

\subsection{从自由采样到偶句异字同韵}

自由模式按温度 0.8、top-k 为 50、top-p 为 1 的设置抽样，最多生成 64 个新标记。有限输出空间可以在取样前通过词表掩码约束。设前缀为 $h$，第 $t$ 步允许集合为 $\mathcal A_t(h)$，则将允许集合外的 logits 设为 $-\infty$，再对允许集合抽样。固定句式在每五个汉字后生成规定标点，末句标点后只允许结束符。于是读者看到的原文保持四句和二十个汉字。

押韵模式用固定的平水韵表建立韵脚搭档。第二句末在只有一个可确认平声韵部且有搭档的字中抽样，第四句末只允许同部的其他字。多读、多部的韵脚在自动诊断中保留为待判。要求全部关键词时，采样为各词选择不同诗句及句内落点，并让偶句保留韵脚位。模型可在预设落点之前自然生成目标词。默认比较四个候选，以原模型给完整输出的平均对数似然减去整句和双字重复惩罚进行排序。以上机制只规定偶句同韵和所需用词，首句入韵和逐位平仄通过后续诊断评估。

\subsection{冻结主测试与规则归因}

主测试在 50 道预先登记的新题上运行，词数为一、二、三的题目分别有 16、17、17 道。每题使用 11、22、33 三个采样种子，四种模式使用同一权重。统计四句五字、偶句异字同韵、原词全覆盖及三项同时达标。原词匹配要求输入词完整出现在同一句，便于检查生成约束。语义切题另用量表。独立目录中的复测对照了全部 600 条输出，原文和自动判定均与封存版一致。

\reportfigure{rules.pdf}{四种解码方式的规则对照。每种模式计 150 条。只限字数及押韵时全词覆盖与自由模式相同，加入关键词约束后才出现明显变化。}{fig:rules}

自由生成已有 148/150 首符合四句五字，92/150 首确认偶句异字同韵，三项同时达标 17/150。只限长度后，四句五字升到 150/150，三项联合仍是 17/150。加押韵后，偶句异字同韵升到 150/150，三项联合到 25/150。再加入用词约束及四候选排序，四项均为 150/150。自由模式按词数分层，一词为 25/48 全词命中，两词、三词各 0/51。默认模式分别达到 48/48、51/51、51/51。这一差别指向模型主动用全多个词的能力，而非句长与韵脚规则已经达到的上限。

在默认模式的 150 首中，完整生成文本互异，覆盖 24 个可确认平声韵部。按冻结训练集查询整首、长诗连续片段及近重复，三项均为 0/150。对严格逐字平仄常格及同格押韵另行核验时，1 首合、146 首不合、3 首待判。此口径使用四套固定模板，历史变格、多音字和首句入韵存在需据诗句判断的情形。八首已归档的独立 AI 格律复核供个案核对。发布的文字标签使用「五言四句 · 偶句押韵」，并保留平仄诊断结果。

\subsection{两份独立 AI 盲评}

最新 92M 的主观补测从冻结题抽取 30 道，每种词数 10 道，种子固定为 11。每题给出自由与默认输出，两位 AI 评阅者各读 60 首且只接收本人的匿名表与 v2 量表。题序和组内次序独立随机化，同题比较保留配对关系。早期 14M 归档包含一份真人 v1 评分和一份同版 AI 参照，其模型、量表和输出均有独立记录。表~\ref{tab:rubricv2} 是本轮使用的指标。

\begin{table}[!htbp]
\centering
\caption{主观评分定义。分值 1 至 5，NA 表示该项不适用。}
\label{tab:rubricv2}
\begin{tabularx}{\linewidth}{@{}p{.17\linewidth}Y p{.24\linewidth}@{}}
\toprule
维度 & 判断依据 & 适用范围\\
\midrule
切题 & 是否在语义上回应题意，接纳同义、转述和可由诗句推得的暗示。 & 全部输出\\
可读性 & 读者能否理解用词和句意，接纳合理的古汉语省略、倒装。 & 全部输出\\
连贯性 & 四句场景、时间、叙述视角或情绪是否能衔接。 & 可解析的四句\\
意境 & 画面、情意和余味是否由诗句自然形成。 & 可读性和连贯性均至少 3 分\\
\bottomrule
\end{tabularx}
\end{table}

\reportfigure{blind.pdf}{两份 AI v2 盲评的配对均差。括号为每位评阅者的可比题数。默认模式的两个切题判断方向一致，可读性有分歧。}{fig:blind}

两个评阅者对默认模式的切题配对均差分别为 $+1.67$ 和 $+1.50$ 分，可读性为 $+0.17$ 与 $-0.07$ 分，连贯性为 $+0.37$ 与 $+0.23$ 分，三维的配对数均为 30。意境的配对均差为 $+0.20$ 与 $+0.31$，两个分母分别为 25、26；这是两种输出均通过适用门槛的子集。两份切题评分的增长提供了语义层面的补充证据，贴近题词和自然表达的关系仍须通过实例及新的真人盲评判断。两个 AI 上下文可能共享判断偏好，分值用于描述这批输出。

\begin{poembox}
\small\noindent\textbf{同题对读}\quad 题词「{\versefont 孤舟、春雨、煙波}」，种子 11\par
\noindent 自由生成\quad {\versefont 春潮平地脈，曉氣入雲深。獨立蒼茫裏，無人識此心。}\par
\noindent 默认模式\quad {\versefont 煙波春渺渺，寂寂孤舟閑。春雨夜來雨，空江寒葦山。}
\end{poembox}

这组输出说明两个指标各有独立用途。自由模式用春潮、晓气和独立远眺写出了季候与心境，三个原词均未完整出现。默认模式为读者提供明确的烟波与孤舟，三个原词均可定位，「春雨夜來雨」读来较生硬。只统计原词无法解释这种分歧。系统需要更多让多个词自然共现的训练配对，以及能在新题上检验衔接和表达的评分。

\subsection{新实现与新题上的隔位重字实验}

原排序已对整句及重复双字扣分，仍可能留下同一句内隔着其他字再次出现的单字。新排序把每句相邻同字合并成一个连续段，再统计该字曾出现过多少次。因而「渺渺」保留为一个连续段，「春雨夜來雨」计一次隔位返回。设原分数为 $s(y)$，计数为 $r(y)$，新排序使用 $s(y)-0.18\,r(y)$。权重在新题生成前固定。两种排序按相同约束、候选布局和随机种子重建同一候选池，并只改变排序规则。CLI 与本地网页均可切换。

实际补测分为两部分。先在已有 90 道验证题（种子 11）检查排序，原排序与新排序的重复分别为 12/90 和 7/90。继而冻结 36 道新题，一、二、三词各 12 道，所有词均通过字表检查，词集合和原主测试互异。每题各运行三个种子。四候选原排序和新排序有 4 条输出变化，新题中的隔位重字由 13/108 降到 9/108。按题作为抽样单位的配对 bootstrap，率差为 $-3.70$ 个百分点，95\% 分位区间为 $[-7.41,-0.93]$ 个百分点。该结果对应被直接优化的重字指标。

\reportfigure{ranking.pdf}{补充实验。A 固定候选后比较两种排序的隔位重字率。B 为新题 GPU 热推理延迟，每个模式 108 次请求。每条线连接中位数和第 95 百分位数。}{fig:ranking}

新题上单候选、四候选原排序和四候选新排序均是 108/108 三项达标。RTX 4080 SUPER 的中位耗时分别为 0.059、0.255、0.262 秒，同句隔位重字分别为 15/108、13/108、9/108。单候选适合更重视等待时间的演示。两种四候选模式的时延接近。自由生成在两词、三词各 36 次仍无全词覆盖。补测支持一个小范围的工程结论，即在现有候选中可少选几首隔位重字明显的诗。个别改善后文本依然会有生硬词组，且本轮新的 AI 盲评用的是原默认排序。新排序的语义与意境仍待同题盲评。
